% !TEX root = ../main.tex
\chapter{Conclusioni e Sviluppi Futuri}
\label{chap:conclusions}

\section{Risultati ottenuti}
Il presente lavoro di tesi ha permesso di esplorare in modo sistematico il profondo legame sussistente tra l'informatica teorica e la logica matematica, 
culminando nella formalizzazione parziale dell'isomorfismo di Curry-Howard all'interno dell'assistente alla dimostrazione Lean 4. Partendo dalle fondamenta 
del $\lambda$-calcolo e dei sistemi di tipi, si è analizzata l'evoluzione sintattica e semantica che conduce al Calcolo delle Costruzioni Induttive (CiC). 
Dal punto di vista pratico, l'indagine ha dimostrato l'efficacia di Lean 4 come ambiente unificato: il framework implementato mappa con successo le formule 
logiche della logica proposizionale intuizionista (IPC) sui tipi ed esprime il processo di semplificazione delle prove come $\beta$-riduzione computazionale. 
Nonostante la deliberata scelta metodologica di lasciare alcune dimostrazioni algebricamente onerose postulate tramite l'operatore \texttt{sorry}, l'architettura 
delle segnature e la simmetria dei tipi dipendenti descritte nel caso di studio confermano la solidità del modello teorico proposto.

\section{Sviluppi futuri}
L'analisi condotta in questa trattazione apre la strada a progetti future focalizzate sulle capacità di Lean 4 come linguaggio di programmazione e assistente di 
dimostrazione per l'automazione dell'analisi del software. In particolare, per stringenti ragioni di tempo, non è stato possibile implementare in questa tesi un 
progetto nativo in Lean 4, volto a sfruttare le strutture algebriche e coalgebriche per il calcolo automatico e formalmente verificato della complessità asintotica 
dei programmi. Di seguito viene delineata l'architettura formale di questo sistema, concepito per unire strumenti analitici e metaprogrammazione.

\subsection{Inferenza automatica della complessità asintotica}
L'obiettivo centrale della componente software pianificata consiste nell'utilizzare l'espressività di Lean 4 per realizzare un sistema di inferenza automatica della complessità 
asintotica delle funzioni. Al fine di garantire l'attendibilità e il rispetto dei principi della verifica formale, il problema viene strutturato in due macro-fasi distinte:

\begin{enumerate}
    \item \textbf{Composizione dei costi tramite monadi di complessità:} 
        La prima parte del sistema affronta la combinazione di funzioni sequenziali. Il costo computazionale non viene calcolato empiricamente, ma viene 
        internalizzato direttamente all'interno degli effetti collaterali del tipo di ritorno mediante l'uso di una specifica \emph{monade di complessità}. 
        Questa struttura monadica si occupa di accumulare e propagare la complessità asintotica man mano che le funzioni vengono composte, trasformando 
        il calcolo del costo in un'operazione guidata dal sistema dei tipi.
    
    \item \textbf{Analisi della ricorsione tramite metaprogrammazione:} 
        La seconda parte del problema richiede l'uso della metaprogrammazione nativa di Lean 4. Attraverso la scrittura di macro e tattiche personalizzate, 
        il software ispeziona l'albero di sintassi astratta (AST) del codice a tempo di compilazione per ``leggere'' la forma della ricorsione. Il meta-programma 
        ha il compito di isolare autonomamente i casi base e i passi ricorsivi, traducendo la struttura sintattica della funzione in un'equazione di ricorrenza formale.
\end{enumerate}
Una volta estratta l'equazione di ricorrenza, il sistema utilizza teoremi e strumenti analitici pre-dimostrati (come le variazioni del \textit{Master Theorem}) per ricavare 
in modo autonomo e matematicamente certo la complessità asintotica finale. L'intero processo garantisce un'affidabilità totale: ogni inferenza prodotta non è una mera stima, 
ma genera un termine di prova che certifica la correttezza del calcolo asintotico rispetto al kernel.

\subsection{Integrazione con Mathlib}
Per mantenere il sistema focalizzato sugli aspetti concettuali e di automazione, l'architettura del progetto rigetta programmaticamente la scrittura di codice non didattico 
o ridondante. Invece di reimplementare i numeri reali ($\mathbb{R}$), le relazioni di disuguaglianza asintotica (la notazione $O$-grande) e il calcolo delle funzioni reali, 
il software delega interamente queste nozioni a \texttt{Mathlib}, la libreria standard di Lean 4. Sfruttando i teoremi dell'analisi già formalizzati e verificati all'interno 
dell'ecosistema di \texttt{Mathlib}, lo sforzo implementativo si concentra esclusivamente sulla logica di metaprogrammazione per l'estrazione delle ricorrenze e sulla loro 
risoluzione automatica.

\subsection{Integrazione nella rappresentazione LCNF}
Infine, per massimizzare l'attendibilità dell'inferenza rispetto all'eseguibile finale, l'analisi sintattica della complessità potrebbe essere spostata dal Frontend 
direttamente alla forma normale del compilatore (Lean Compiler Normal Form - LCNF), analizzata nel Capitolo \ref{chap:architecture}. In questo modo, l'inferenza della 
complessità terrebbe conto delle ottimizzazioni reali introdotte dal compilatore (come l'inlining e il \textit{destructive update} tramite reference counting), valutando 
il comportamento effettivo a runtime e non la mera astrazione del codice sorgente inserito dall'utente.